% ===== 2 总体方案与技术路线 =====
\section{总体方案与技术路线}
\label{sec:system-design}

\subsection{技术路线比较与选择}
\label{sec:route-selection}

面向本赛题，技术路线的选择不能只依据单一图像分割指标，而需要同时考虑自然堆积颗粒的接触与遮挡、比赛现场可获得的数据量、标准筛分的物理评价口径以及有限的部署时间。综合现有研究与工程实现方式，可将候选方案概括为传统图像处理、监督式实例分割、预训练通用/颗粒分割模型以及多视角或三维测量四类。表\ref{tab:route-comparison}从与本赛题直接相关的维度比较了这些路线。

\begin{table}[htbp]
\centering
\caption{面向竞赛约束的候选技术路线比较}
\label{tab:route-comparison}
\footnotesize
\renewcommand{\arraystretch}{1.15}
\begin{tabular}{>{\raggedright\arraybackslash}p{2.7cm}>{\raggedright\arraybackslash}p{2.8cm}>{\raggedright\arraybackslash}p{2.6cm}>{\raggedright\arraybackslash}p{2.4cm}>{\raggedright\arraybackslash}p{3.2cm}}
\toprule
技术路线 & 接触颗粒处理 & 数据与训练需求 & 现场复杂度 & 对本赛题的主要权衡 \\
\midrule
阈值/边缘/分水岭 & 对弱接触场景有效，强粘连时易合并或过分割 & 无需训练，但参数依赖场景 & 低 & 实现简单、速度快，但对光照、背景和堆积状态敏感 \\
监督式检测/实例分割 & 在同分布数据上可获得较高精度 & 需要检测框或实例掩膜标注及训练 & 中 & 可针对赛题优化，但现场域变化与标注成本较高 \\
预训练通用/颗粒分割 & 可直接恢复实例轮廓，对跨场景泛化更友好 & 可零样本使用，必要时再微调 & 中低 & 适合快速部署，但下游仍需建立物理尺度和筛分语义 \\
多视角/三维点云 & 可显式恢复三维尺寸与部分厚度信息 & 需额外标定、配准或三维算法 & 高 & 物理表征更完整，但设备、标定与现场时间成本更高 \\
\bottomrule
\end{tabular}
\end{table}

传统阈值、边缘检测和分水岭等方法具有计算开销低、实现透明的优点，但其判别依据主要来自局部灰度、颜色和梯度。当颗粒颜色接近、接触边界较弱或阴影与真实边界混叠时，同一套阈值和形态学参数往往难以跨图像保持稳定。因此，这类方法更适合作为受控成像条件下的轻量基线，而不宜承担自然堆积场景的核心实例感知任务。

监督式目标检测和实例分割方法能够通过数据驱动方式学习颗粒外观，其中 Mask R-CNN 等实例分割模型可以直接预测颗粒级掩膜\cite{he2017maskRCNN}，相关研究也已将深度学习用于骨料和岩块粒度分析\cite{zhang2020aggregate,ferello2023rock,zeng2022automation}。这一路线的主要限制并非模型表达能力不足，而是比赛现场的训练数据和时间条件：若要依靠监督训练获得稳定性能，需要具有代表性的颗粒实例标注，并保证训练数据与现场相机、光照、石料类型和堆积方式具有足够一致性。对于强调快速部署的竞赛场景，这会引入额外的数据准备风险。

近年来，通用分割模型和面向颗粒的预训练模型显著降低了重新标注和训练的门槛。Segment Anything 展示了大规模预训练分割模型的通用目标表征能力\cite{kirillov2023segmentAnything}；Cellpose 系列则强调对形态多样实例的泛化分割\cite{stringer2021cellpose,pachitariu2025cellposeSAM}。ImageGrains 2.0 在跨类型颗粒数据上构建了面向粒状材料的分割与测量流程，并提供基于 Cellpose-SAM 的预训练模型\cite{mair2026imagegrains2,zenodoIG2}。因此，本文选择 ImageGrains 2.0 / Cellpose-SAM 作为实例感知基座：其目的不是假设预训练模型已经解决全部筛分问题，而是优先利用已有颗粒分割能力解决最依赖大规模视觉数据的边界识别问题，将有限的赛题数据集中用于后续物理校准和精度验证。

另一类路线通过多视角融合或三维点云直接估计颗粒三维几何，可在原理上减少单目投影造成的厚度信息缺失\cite{fan2021multiview,yang2023pointCloud}。然而，多相机/深度设备的外参标定、点云配准与现场布置会显著增加系统复杂度。考虑本赛题的现场准备时间和系统可重复性，本文首先采用结构更简单的单目俯视方案，并通过显式筛分等效模型吸收二维投影到机械筛分之间的系统偏差；三维测量则作为后续提高物理真实性的重要扩展方向。

综上，本文的路线选择遵循“将强数据需求留给预训练视觉模型，将赛题特有关系留给显式校准”的思路：上游采用具有跨场景泛化能力的颗粒实例分割模型，下游不直接将视觉尺寸视为标准筛分结果，而是通过可解释的尺度恢复、筛分等效粒径和质量加权模型完成工程语义转换。完成核心技术路线选择后，还需要进一步明确各模块在系统层面共同遵循的设计原则。

\subsection{系统设计原则}
\label{sec:design-principles}

本系统的最终目标并不是在某一固定测试图像上获得尽可能复杂的预测结果，而是在未知现场样本上以可复现方式输出能够与标准筛分比较的结果。因此，整体设计由\textbf{泛化优先、校准显式和部署可控}三个原则共同约束；三者分别对应上游视觉感知、下游工程映射以及比赛现场执行中的主要风险。

\textbf{泛化优先。} 视觉感知模块应尽量降低对比赛现场重新标注和训练的依赖。自然堆积骨料在岩性、颜色、表面纹理、粒径和照明条件上都可能发生变化，如果核心分割能力依赖少量现场样本重新训练，则有限调试时间会同时承担数据采集、标注、训练和模型选择等不确定性。本文因此优先使用已经在多类型颗粒数据上训练的 ImageGrains 2.0 / Cellpose-SAM，仅保留少量与成像尺度和实例后处理有关的可调参数，并将模型输出限制为后续可检查的实例掩膜，而不是直接产生不可追溯的最终级配结论。

\textbf{校准显式。} 凡是直接决定毫米尺度或标准筛分质量口径的转换关系，均采用低维参数化模型并保留真实物理试验的校准入口。像素到毫米的转换由尺度系数明确控制；二维颗粒几何到筛分等效粒径由参数 $\boldsymbol{\theta}$ 控制；颗粒尺寸到质量代理权重由指数 $\gamma$ 控制。与将图像直接输入端到端网络预测级配曲线的方式相比\cite{coenen2022learningToSieve,buscombe2020sedinet}，这种设计牺牲了一部分模型自由度，但能够明确回答误差来自“视觉边界”“尺度”“粒径映射”还是“质量假设”，也便于在获得少量标准筛分批次后重新拟合关键参数。已有数字图像粒度研究同样表明，从数量型图像统计到重量型粒级分布需要额外的质量转换模型\cite{maiti2017massModel}。

\textbf{部署可控。} 系统应在满足准确性的前提下尽量压缩现场硬件、操作和参数搜索复杂度。本文采用单目俯视成像、固定或可快速恢复的成像几何、统一尺度参考物以及模块化命令行处理流程，将现场步骤限制为设备布置、尺度校验、图像采集、结果质检和自动分析几个可重复环节。算法性能也不只以神经网络前向推理时间衡量，而是在后续实验中统计从现场准备到报告输出的完整 wall-clock time，以直接对应比赛规定的时间窗口。

上述三个原则决定了本文不采用“单一端到端网络解决全部任务”的结构，而是按照数据依赖和误差来源将系统拆分为若干可检查、可替换的模块。下一节给出从现场图像到最终报告的完整数据流，并说明标准筛分真值如何作为离线反馈重新进入该流程。

\subsection{整体系统架构}
\label{sec:overall-architecture}

本文系统由图像采集与尺度恢复、实例感知、颗粒几何测量、筛分级配映射、形貌/异常分析以及结果汇总六个功能阶段组成。其主数据流是一条从原始图像到场景级筛分报告的单向处理链，而机械筛分真值不参与单次在线推理，仅在离线阶段用于校准筛分等效粒径和质量权重参数。图\ref{fig:overall-architecture}给出了整体架构及两类数据流之间的关系。

\begin{figure}[htbp]
\centering
\begin{tikzpicture}[
  node distance=0.65cm and 0.55cm,
  block/.style={rectangle, draw, rounded corners, minimum height=1.0cm, minimum width=2.55cm, align=center, font=\small},
  smallblock/.style={rectangle, draw, rounded corners, minimum height=0.8cm, minimum width=2.3cm, align=center, font=\scriptsize},
  arrow/.style={->, thick},
  feedback/.style={->, thick, dashed}
]
\node[block, fill=blue!7] (acq) {图像采集\\与尺度参考};
\node[block, fill=blue!7, right=of acq] (seg) {实例分割\\$I\rightarrow\{M_i\}$};
\node[block, fill=green!7, right=of seg] (geo) {物理几何测量\\$M_i\rightarrow\mathbf{x}_i$};
\node[block, fill=orange!8, right=of geo] (sieve) {筛分语义映射\\$\mathbf{x}_i\rightarrow d_i,w_i$};

\node[smallblock, fill=orange!8, below=0.8cm of sieve] (grading) {质量级配\\$D_p,\;P_k$};
\node[smallblock, fill=purple!7, left=0.45cm of grading] (shape) {形貌分类\\$c_i$};
\node[smallblock, fill=red!6, right=0.45cm of grading] (anomaly) {异常检测\\$z_i$};
\node[block, fill=gray!8, below=0.85cm of grading] (report) {场景级结果汇总\\表格、曲线与报告};

\draw[arrow] (acq) -- node[above, font=\scriptsize]{图像+尺度} (seg);
\draw[arrow] (seg) -- node[above, font=\scriptsize]{实例掩膜} (geo);
\draw[arrow] (geo) -- node[above, font=\scriptsize]{毫米级特征} (sieve);
\draw[arrow] (sieve) -- (grading);
\draw[arrow] (geo.south) |- (shape.west);
\draw[arrow] (geo.south) |- (anomaly.west);
\draw[arrow] (shape) |- (report.west);
\draw[arrow] (grading) -- (report);
\draw[arrow] (anomaly) |- (report.east);

\node[smallblock, fill=yellow!12, below=0.75cm of acq] (gt) {同批次标准\\机械筛分真值};
\node[smallblock, fill=yellow!12, below=0.75cm of seg] (calib) {离线参数校准\\$\boldsymbol{\theta},\gamma$};
\draw[feedback] (gt) -- (calib);
\draw[feedback] (calib.east) -| (sieve.south west);
\end{tikzpicture}
\caption{总体系统架构。实线表示单次在线分析的数据流；虚线表示利用同批次机械筛分真值进行的离线参数校准。机械筛分真值不作为上游分割网络输入，而只用于校准视觉几何到质量级配之间的工程映射。}
\label{fig:overall-architecture}
\end{figure}

在主处理链的前半部分，现场图像首先通过已知尺寸参考物建立像素到毫米的尺度映射，随后由 Cellpose-SAM 恢复每个可见颗粒的实例掩膜。实例掩膜进一步被转换为统一的几何描述，包括长轴、短轴、投影面积、周长、圆度和凸度等。这样，上游视觉模块的输出不是某一个最终类别或粒径，而是一组可以被人工检查、可以被不同下游任务复用的颗粒级中间表示。

在后半部分，统一几何表征被分为两类处理路径。主路径面向比赛核心粒径任务：通过筛分等效粒径模型得到 $d_i$，再通过质量代理权重 $w_i$ 将颗粒级预测聚合为累计质量级配以及各标准筛级占比。并行路径则复用同一套几何特征完成形貌分类和异常筛查，避免为每个附加任务重新建立独立的检测网络。所有任务结果最终汇总到场景级报告中，并同时保留实例掩膜、颗粒表格和中间可视化，以便现场快速核验。

标准机械筛分构成一条与在线推理相互独立的离线反馈链。对于若干具有配对真值的骨料批次，机械筛分给出真实累计质量分布，视觉系统给出相应颗粒级几何统计；二者之间的误差用于优化 $\boldsymbol{\theta}$ 和 $\gamma$。这种设计有两个直接好处：其一，新增筛分真值时无需重新训练计算量更大的实例分割网络；其二，可以在独立测试批次上严格区分“参数校准”和“性能验证”，避免使用同一批数据既拟合又评价最终级配精度。

因此，整个系统可以理解为一个由\textbf{通用视觉感知}和\textbf{赛题特定物理统计校准}组成的组合架构。前者负责尽可能稳定地回答“每颗颗粒在哪里、轮廓是什么”，后者负责回答“这些轮廓在标准筛分意义下代表什么”。为了使这一架构与评审任务形成明确闭环，下一节进一步给出各系统模块的输入、输出以及对应验证方式。

\subsection{赛题任务与系统模块对应关系}
\label{sec:module-task-mapping}

第\ref{sec:task-background}节已经从比赛要求出发说明了系统需要输出哪些结果；本节进一步从系统实现角度规定每个核心模块应接受什么输入、产生什么输出，以及最终通过何种独立证据验证。这样可以将“功能已经实现”和“功能已经被证明有效”明确区分：可运行的软件输出只证明处理链闭环，而识别准确性和筛分准确性必须由相应人工标注或物理试验真值进行评价。

\begin{table}[htbp]
\centering
\caption{系统模块、输出与验证证据的对应关系}
\label{tab:module-validation}
\footnotesize
\renewcommand{\arraystretch}{1.18}
\begin{tabular}{>{\raggedright\arraybackslash}p{2.5cm}>{\raggedright\arraybackslash}p{3.0cm}>{\raggedright\arraybackslash}p{3.6cm}>{\raggedright\arraybackslash}p{4.6cm}}
\toprule
系统模块 & 主要输入 & 主要输出 & 后续验证依据 \\
\midrule
实例感知 & 原始图像 $I$ & 实例掩膜 $\{M_i\}$ & 人工实例标注；IoU、Precision/Recall、merge/split 错误 \\
尺度与几何测量 & 掩膜、尺度参考 & $a_i,b_i,A_i,P_i,C_i,S_i$ & 已知长度参考物、部分颗粒人工尺寸测量 \\
筛分语义映射 & 颗粒几何、$\boldsymbol{\theta},\gamma$ & $d_i,w_i,D_p,P_k$ & 同批次标准机械筛分；$D_p$ 与筛级质量占比误差 \\
形貌分类 & 颗粒几何特征 & 形貌类别 $c_i$ & 人工/专家形貌标签；Macro-F1、混淆矩阵 \\
异常检测 & 粒径与形貌特征/颗粒裁剪 & 异常状态 $z_i$ & 已知异常样本；Precision、Recall、F1 \\
工程工作流 & 软硬件与现场样本 & 图表、报告、完整处理结果 & 从设备布置到报告输出的 wall-clock time 与重复运行稳定性 \\
\bottomrule
\end{tabular}
\end{table}

表\ref{tab:module-validation}同时定义了本文后续实验章节的证据结构。实例分割是否“看起来合理”不能替代人工掩膜评价，预测级配是否“趋势合理”不能替代标准机械筛分对照；同理，当前基于几何规则得到的形貌和疑似泥团标签也只有在获得独立人工标注后才能报告严格分类精度。本文将按照这一原则分别区分功能演示、参数校准与独立精度验证，避免将模型内部派生量作为自身 ground truth。

至此，系统层面的技术选择、设计原则、数据流和验证闭环均已明确。沿图\ref{fig:overall-architecture}的主数据流，下一章首先讨论如何通过成像控制、尺度标定、实例分割和几何特征提取，将现场照片稳定转换为毫米尺度下的颗粒级物理表征。